\DocumentMetadata{
  pdfversion=2.0,
  pdfstandard=ua-2,
  lang=en-US,
  tagging=on,
  tagging-setup={math/setup=mathml-SE,tabsorder=structure}
}
\documentclass[11pt,letterpaper]{article}
\usepackage[margin=0.68in,includefoot]{geometry}
\usepackage{newcomputermodern}
\usepackage{amsmath,amscd,mathtools}
\usepackage{tikz,adjustbox}
\providecommand{\square}{\mdlgwhtsquare}
\usepackage[hidelinks]{hyperref}
\hypersetup{
  pdftitle={Algebra First-Year Exam, May 2025, Part 2},
  pdfauthor={University of Florida Department of Mathematics},
  pdfsubject={Graduate mathematics examination},
  pdfdisplaydoctitle=true,
  bookmarksopen=true
}
\setcounter{secnumdepth}{0}
\setlength{\parindent}{0pt}
\setlength{\parskip}{0.28em}
\setlength{\emergencystretch}{3em}
\setlength{\leftmargini}{2.15em}
\setlength{\leftmarginii}{2.65em}
\setlength{\leftmarginiii}{3em}
\pagestyle{empty}
\raggedbottom
\allowdisplaybreaks
\NewTaggingSocketPlug{block/list/label}{examproblem}{%
  \tagstructbegin{tag=Lbl}\tagstructend%
  \tagstructbegin{tag=\LBody}%
  \tagstructbegin{tag=\ProblemHeadingTag,title-o={\ProblemHeadingText}}%
  \tagmcbegin{tag=\ProblemHeadingTag,actualtext={\ProblemHeadingText}}%
  #2%
  \tagmcend%
  \tagstructend%
}
\newenvironment{examproblems}[2]{%
  \def\ProblemHeadingTag{#2}%
  \AssignTaggingSocketPlug{block/list/label}{examproblem}%
  \begin{enumerate}%
  \setcounter{enumi}{#1}%
  \renewcommand{\labelenumi}{\textbf{\theenumi.}}%
  \setlength{\topsep}{0.35em}%
  \setlength{\partopsep}{0pt}%
  \setlength{\itemsep}{0.48em}%
  \setlength{\parsep}{0.18em}%
}{\end{enumerate}}
\newenvironment{examsubparts}{%
  \AssignTaggingSocketPlug{block/list/label}{default}%
  \begin{enumerate}%
  \setlength{\topsep}{0.18em}%
  \setlength{\partopsep}{0pt}%
  \setlength{\itemsep}{0.16em}%
  \setlength{\parsep}{0.1em}%
}{\end{enumerate}}
\newenvironment{examinstructions}{%
  \par\begingroup\itshape
}{%
  \par\endgroup\vspace{0.18em}
}
\makeatletter
\renewcommand\subsection{\@startsection{subsection}{2}{0pt}%
  {-1.25ex plus -.25ex minus -.1ex}%
  {0.55ex plus .1ex}%
  {\normalfont\large\bfseries}}
\renewcommand{\@maketitle}{%
  \newpage\null\vskip -1em%
  {\footnotesize\bfseries
    \href{https://www.ufl.edu/}{University of Florida}, \href{https://math.ufl.edu/}{Department of Mathematics}\par}%
  \vskip -0.5em%
  \tagstructbegin{tag=H1,title={Algebra First-Year Exam, May 2025, Part 2}}%
  \tagpdfparaOff%
  \tagmcbegin{tag=H1}%
    {\LARGE\bfseries\href{https://gma.math.ufl.edu/exams/algebra-first-year/}{Algebra First-Year Exam}, May 2025, Part 2\par}%
  \tagmcend%
  \tagpdfparaOn%
  \tagstructend%
  \vskip 0.25em%
  \tagmcbegin{artifact}%
    \hrule height 1pt%
  \tagmcend%
  \vskip 1em%
}
\makeatother
\title{Algebra First-Year Exam, May 2025, Part 2}
\author{}
\date{}
\begin{document}
\maketitle
\thispagestyle{empty}
\begin{examinstructions}
Answer 4 questions. Write your solutions clearly in complete English sentences. You may quote theorems (within reason) as long as you state them clearly.
\end{examinstructions}
\begin{examproblems}{0}{H2}
\def\ProblemHeadingText{Problem 1.}
\item
This problem has two parts.
\begin{examsubparts}
\item[\textnormal{(a)}]
Define the terms Euclidean domain and principal ideal.
\item[\textnormal{(b)}]
Prove that in a Euclidean domain every ideal is a principal ideal.
\end{examsubparts}
\def\ProblemHeadingText{Problem 2.}
\item
Let~\(p\) be a prime and \(R=\{r\in\mathbb{Q}\mid r=\frac{a}{b}\text{ for some }a,b\in\mathbb{Z}\text{ with }p\nmid b\}\). Prove that \(pR\) is the unique maximal ideal of~\(R\).
\def\ProblemHeadingText{Problem 3.}
\item
Using the rational canonical form, find one representative for each conjugacy class in \(\operatorname{GL}(4,\mathbb{F}_2)\) whose elements have order~\(3\) or~\(6\). (Hint: in \(\mathbb{F}_2[x]\) we have \(x^6-1=(x^3-1)^2\).)
\def\ProblemHeadingText{Problem 4.}
\item
Determine which of the following are irreducible in the given rings.
\begin{examsubparts}
\item[\textnormal{(a)}]
\(X^3-X^2+X+2\) in \(\mathbb{Q}[X]\)
\item[\textnormal{(b)}]
\(X^4+X+1\) in \(\mathbb{F}_2[X]\)
\item[\textnormal{(c)}]
\(Y^3+2X^2Y+X(X+1)\) in \(\mathbb{Q}[X,Y]\)
\end{examsubparts}
\def\ProblemHeadingText{Problem 5.}
\item
Let~\(F\) be a field, \(V\) a finite-dimensional~\(F\)-vector space and \(T\in\operatorname{End}_F(V)\) a linear operator.
\begin{examsubparts}
\item[\textnormal{(a)}]
Explain why there is an~\(F\)-algebra homomorphism \(\psi_T:F[x]\to\operatorname{End}_F(V)\) mapping~\(x\) to~\(T\).
\item[\textnormal{(b)}]
Define the minimum polynomial \(m_T(x)\) of~\(T\).
\item[\textnormal{(c)}]
Let \(F[T]\) be the~\(F\)-subalgebra of \(\operatorname{End}_F(V)\) generated by~\(T\). Prove that if \(m_T(x)\) is irreducible in \(F[x]\) then every nonzero element of \(F[T]\) is invertible (under composition of operators).
\end{examsubparts}
\def\ProblemHeadingText{Problem 6.}
\item
Let \(n\geq 1\) and let \(F=\mathbb{Q}(\sqrt{1},\sqrt{2},\ldots,\sqrt{n})\). Show that \(\sqrt[3]{2}\notin F\).
\end{examproblems}
\end{document}
